
\documentclass{article}
\usepackage{colortbl}
\usepackage{makecell}
\usepackage{multirow}
\usepackage{supertabular}

\begin{document}

\newcounter{utterance}

\centering \large Interaction Transcript for game `cladder', experiment `full\_v1.5\_default', episode 1169 with llama{-}3.1{-}8b{-}instruct.
\vspace{24pt}

{ \footnotesize  \setcounter{utterance}{1}
\setlength{\tabcolsep}{0pt}
\begin{supertabular}{c@{$\;$}|p{.15\linewidth}@{}p{.15\linewidth}p{.15\linewidth}p{.15\linewidth}p{.15\linewidth}p{.15\linewidth}}
    \# & \multicolumn{2}{c}{Player} && \multicolumn{2}{c}{Game Master} \\
    \hline

    \theutterance \stepcounter{utterance}  
    & & & \multicolumn{4}{p{0.6\linewidth}}{
        \cellcolor[rgb]{0.9,0.9,0.9}{
            \makecell[{{p{\linewidth}}}]{
                \texttt{\tiny{[P1$\langle$GM]}}
                \texttt{You are an expert in causal inference. The following question is not a typical commonsense query, but rather a meticulously designed question created by a professor specializing in causal inference, intended to assess the students' mastery of the course content.} \\
\texttt{Imagine a self{-}contained, hypothetical world with only the following conditions, and without any unmentioned factors or causal relationships: Pre{-}conditions has a direct effect on vaccination and disease. Vaccination has a direct effect on disease. The overall probability of getting the vaccine is 59\%. The probability of vaccine refusal and recovering from the disease is 24\%. The probability of getting the vaccine and recovering from the disease is 20\%. Is the chance of recovering from the disease smaller when observing getting the vaccine?} \\
\texttt{Start your answer with 'Yes' or 'No', followed by additional reasoning or evidence to support your explanation.} \\
            }
        }
    }
     \\ \\

    \theutterance \stepcounter{utterance}  
    & \multicolumn{4}{p{0.6\linewidth}}{
        \cellcolor[rgb]{0.561,0.737,0.561}{
            \makecell[{{p{\linewidth}}}]{
                \texttt{\tiny{[P1$\rangle$GM]}}
                \texttt{No.} \\
\\ 
\texttt{To assess whether the chance of recovering from the disease is smaller when observing getting the vaccine, we need to apply the concept of conditional probability and potential outcomes. } \\
\\ 
\texttt{Let's denote the event of getting the vaccine as V, the event of recovering from the disease as R, and the event of not getting the vaccine as NV. We are given the following probabilities:} \\
\\ 
\texttt{{-} P(V) = 0.59 (probability of getting the vaccine)} \\
\texttt{{-} P(NV) = 1 {-} P(V) = 0.41 (probability of not getting the vaccine)} \\
\texttt{{-} P(R|NV) = 0.24 (probability of recovering from the disease given no vaccination)} \\
\texttt{{-} P(R|V) = 0.20 (probability of recovering from the disease given vaccination)} \\
\texttt{{-} P(V ∩ R) = 0.20 (probability of getting the vaccine and recovering from the disease)} \\
\\ 
\texttt{We can calculate the probability of recovering from the disease given vaccination (P(R|V)) and the probability of recovering from the disease given no vaccination (P(R|NV)).} \\
\\ 
\texttt{However, to determine whether the chance of recovering from the disease is smaller when observing getting the vaccine, we need to compare the probability of recovering from the disease given vaccination (P(R|V)) with the probability of recovering from the disease given no vaccination (P(R|NV)).} \\
\\ 
\texttt{To do this, we can use the law of total probability to calculate the overall probability of recovering from the disease (P(R)):} \\
\\ 
\texttt{P(R) = P(R|V) * P(V) + P(R|NV) * P(NV)} \\
\texttt{= 0.20 * 0.59 + 0.24 * 0.41} \\
\texttt{= 0.118 + 0.0984} \\
\texttt{= 0.2164} \\
\\ 
\texttt{Now, we can calculate the probability of recovering from the disease given no vaccination (P(R|NV)):} \\
\\ 
\texttt{P(R|NV) = P(R ∩ NV) / P(NV)} \\
\texttt{= (P(R|NV) * P(NV)) / P(NV)} \\
\texttt{= 0.24 * 0.41 / 0.41} \\
\texttt{= 0.24} \\
\\ 
\texttt{We can also calculate the probability of recovering from the disease given vaccination (P(R|V)):} \\
\\ 
\texttt{P(R|V) = P(R ∩ V) / P(V)} \\
\texttt{= (P(R|V) * P(V)) / P(V)} \\
\texttt{= 0.20 * 0.59 / 0.59} \\
\texttt{= 0.20} \\
\\ 
\texttt{Comparing P(R|V) and P(R|NV), we can see that P(R|V) = 0.20 and P(R|NV) = 0.24. Since P(R|V) < P(R|NV), the chance of recovering from the disease is actually smaller when observing getting the vaccine.} \\
\\ 
\texttt{However, the question asks whether the chance of recovering from the disease is smaller when} \\
            }
        }
    }
    & & \\ \\

    \theutterance \stepcounter{utterance}  
    & & & \multicolumn{2}{p{0.3\linewidth}}{
        \cellcolor[rgb]{0.9,0.9,0.9}{
            \makecell[{{p{\linewidth}}}]{
                \texttt{\tiny{[GM$|$GM]}}
                \texttt{No} \\
            }
        }
    }
    & & \\ \\

    \theutterance \stepcounter{utterance}  
    & & & \multicolumn{2}{p{0.3\linewidth}}{
        \cellcolor[rgb]{0.9,0.9,0.9}{
            \makecell[{{p{\linewidth}}}]{
                \texttt{\tiny{[GM$|$GM]}}
                \texttt{yes} \\
            }
        }
    }
    & & \\ \\

    \theutterance \stepcounter{utterance}  
    & & & \multicolumn{2}{p{0.3\linewidth}}{
        \cellcolor[rgb]{0.9,0.9,0.9}{
            \makecell[{{p{\linewidth}}}]{
                \texttt{\tiny{[GM$|$GM]}}
                \texttt{game\_result = LOSE} \\
            }
        }
    }
    & & \\ \\

\end{supertabular}
}

\end{document}
